\section{Generating Criticism}
\label{sec:critic}

Steps 10 to 12 of the loop of Section 1 turn from writing to
attack.
The first and most expensive critic is a human: step 10 is an
offline read of 30 to 60 minutes, social media off, and it
cannot be delegated.
The cheapest attacker, used at steps 5 and 11, is an LLM told to
be hostile.
Figure~\ref{fig:critic} shows the prompt: the model plays
Reviewer 2 at the target venue, and the prompt embeds the
venue's own review criteria.
For ICSE'27 those criteria are five: novelty, rigor, relevance,
verifiability and transparency, and presentation; the track's
outcomes are accept, major revision, and reject.
Structuring the critique under the venue's headings matters,
since those headings are exactly the form in which the real
rejection would arrive.
Reviewer 2 is one of four critics the loop assembles; the others
compare the draft against the venue's award winners, the ACM
empirical standards, and the ten nearest prior papers.
The same figure shows that critic running on the generated
drafts of Figures~\ref{fig:sample} and~\ref{fig:samplelit};
each finding quotes the text it fails.

\begin{figure}[t]
\fbox{\parbox{\dimexpr\linewidth-2\fboxsep-2\fboxrule}{\small
{\bfseries Prompt.}
{\ttfamily You are Reviewer 2 for the ICSE'27 research track.
Below are the track's five review criteria, quoted from its call
for papers, and a draft.
Assess the draft under each criterion in turn, quoting the text
that fails it.
End with one of the track's three outcomes: accept, major
revision, or reject.}\\[0.4em]
{\bfseries Output} (on the drafts of Figures~\ref{fig:sample}
and~\ref{fig:samplelit}; abridged):\\[0.2em]
{\bfseries Novelty.}
The FOCUS christening claims a new instrument, yet offers no
delta against existing repertory-grid tooling; the contribution
may be assembly, not invention.\\[0.2em]
{\bfseries Rigor.}
Every headline number is a bracketed placeholder ([17 of 19],
[83\%]): the paper claims an evaluation it has not run, names no
statistical test, and uses one rater.\\[0.2em]
{\bfseries Relevance.}
Matching people to tasks matters to every team; however, evidence
from one six-person team, studying itself, bounds the claimed
impact.\\[0.2em]
{\bfseries Verifiability and transparency.}
Six papers (Giray, Zhang, Xu, Amaral, C\^amara, Khojah) carry no
citations, so none can be verified; RQ2's ground truth, ``the
team's next accepted papers'', cannot exist at submission time.
The promised replication package is the one strong point.\\[0.2em]
{\bfseries Presentation.}
Clear, and the grammar labels aid navigation; minor: the
introduction defines the grid twice.\\[0.2em]
{\bfseries Outcome: reject}, with encouragement to resubmit once
the evaluation exists.}}
\caption{The reviewer-2 critic.
The prompt embeds the five review criteria of the ICSE'27
research track (novelty, rigor, relevance, verifiability and
transparency, presentation), quoted from its call for papers
(conf.researchr.org/track/icse-2027/icse-2027-research-track);
the output, abridged, assesses the step-8 drafts of
Figures~\ref{fig:sample} and~\ref{fig:samplelit} under those
same headings and ends with one of the track's three outcomes.}
\label{fig:critic}
\end{figure}

The review is half the service; triage is the other half.
Figure~\ref{fig:fixes} shows the follow-up prompt.
It divides the review's findings into part 1, fixes the model
can apply alone, and part 2, issues that need a human decision.
The split prices the remaining work.
Clarifications, reorderings, prunings, and citation lookups are
cheap; experiments, ground-truth choices, and scope decisions
are not.
The loop applies part 1 automatically at step 11 and queues
part 2 for the humans at step 12.

\begin{figure}[t]
\fbox{\parbox{\dimexpr\linewidth-2\fboxsep-2\fboxrule}{\small
{\bfseries Prompt.}
{\ttfamily Here is a reviewer's report on the draft.
Whatever its outcome (reject or major revision), divide its
findings into part 1, fixes you can apply yourself
(clarification, reordering, pruning, citation lookup), and
part 2, issues that need a human decision (experiments, ground
truth, scope).
Apply nothing; output the two lists.}\\[0.4em]
{\bfseries Output} (on the review of Figure~\ref{fig:critic};
abridged):\\[0.2em]
{\bfseries Part 1} (the model applies): find and verify
citations for the six named papers; prune the duplicated grid
definition; reorder the RQ2 answer so evidence precedes claim;
use one term, self-report, throughout.\\[0.3em]
{\bfseries Part 2} (the humans decide): run the experiments
that earn the bracketed numbers; pick a defensible ground truth
for RQ2; recruit an independent rater for the constructs; add
more teams, or shrink every generalization claim.}}
\caption{The triage prompt splits the review of
Figure~\ref{fig:critic}: part 1 is applied at step 11 of the
loop; part 2 goes to the humans at step 12.}
\label{fig:fixes}
\end{figure}

Two steps then close the loop.
At step 13, the fixes are applied, and one automatic check
guards the paper's front door: the title and abstract are recoded
with the same flags as the reading set, and they must match the
body's coding with zero flips.
At step 14, a failed self-test returns the loop to step 9;
otherwise the paper ships, with its replication package.
